\documentclass{article}

\usepackage[margin=1in]{geometry}
\usepackage{amsmath}
\usepackage{amssymb}
\usepackage{microtype}

\title{Rewritten Chain of Thought for the\\
Two-Minus Water-Wave Amplitude}
\author{}
\date{}

\newcommand{\pp}[1]{\left[\,#1\,\right]_{+}}

\begin{document}

\maketitle

\begin{abstract}
This document is a first-person reconstruction of the chain of thought for an AI
attempt at the closed-form tree amplitude $A_n$ in the two-minus sector of
one-dimensional deep-water surface waves. It is rebuilt from the visible run
trajectory (the raw private reasoning was encrypted, so only the narrated
messages, tool calls, and outputs survive), and it keeps the original wording,
the false starts, and the hedges rather than tidying them away. The honest
outcome is partial. The prompt's hint was emphatic that $A_n$ is a \emph{single
global rational function} with simple poles on factorization channels and
\emph{no} chamber/piecewise structure, and instructed me to fit one $N/D$ ansatz
on generic kinematics. I followed that route, watched it fail in a specific and
diagnosable way, and concluded the opposite of the hint: the amplitude is
genuinely \emph{piecewise}, controlled by the $|k_S|$ absolute values in the
Berends--Giele propagators. I then proved non-rationality, extracted an exact
general-$n$ closed form valid in the principal (softest-leg-is-a-minus-leg)
regime, and a complete all-chamber formula \emph{only at $n=5$}. The fully
general all-$n$, all-kinematics closed form I did \emph{not} reach; I document
exactly where I stopped and why.
\end{abstract}

\section*{Rewritten Chain of Thought}

I am allowed to read only \texttt{prompt.md} and \texttt{OnShellBG.m}, no sibling
files, no parent files, no web. I read both. The task is a single global
closed-form $A_n$ for $\sigma=(-1,-1,+1,\dots,+1)$, valid for all $n\ge4$ and all
kinematics in the sector. The hint is unusually insistent: $A_n$ is a
\emph{rational function} $N(\omega)/D(\omega)$, the denominator $D$ is the product
of physical factorization-channel factors so $A_n$ carries simple poles on those
channels, the efficient route is one global ansatz fit on a few hundred
\emph{generic} points, and --- stated three times in different words --- there is
``no piecewise / chamber decomposition, no absolute values, and no $\min$/$\max$,''
while the answer is ``emphatically not a plain polynomial.'' I am even told to
\emph{avoid} hierarchical or near-degenerate configurations because they only
ill-condition the fit. I take the hint at face value to begin with; that is the
honest thing to record. It is also, as I will find, wrong, and the way it is
wrong is instructive.

First I check the environment: \texttt{wolframscript}, \texttt{math}, and an
anaconda Python with \texttt{numpy}/\texttt{sympy}/\texttt{mpmath} are present, 64
cores. I re-type the kernel, vertex, propagator, and recursion definitions from
\texttt{OnShellBG.m} into a definitions-only \texttt{bg\_defs.m} (re-typing what I
read is permitted) and run a ground-truth driver over a handful of rational
points. The first surprise is $n=4$: \texttt{BGAmplitude} returns
\texttt{Indeterminate} --- ``\texttt{Power::infy} $1/0$'' and
``\texttt{Infinity::indet} $0\cdot\text{ComplexInfinity}$.'' The two-minus
manifold is rigid at $n=4$: solving the two conservation laws forces the four
frequencies into two opposite pairs $\omega=(-b,a,b,-a)$, so the sub-channels
$\{1,3\}$ and $\{2,4\}$ have $\omega_S=k_S=0$ simultaneously and the propagator
hits a genuine $0/0$. I set $n=4$ aside, to recover later as a limit, and note
that $n=5,6$ give clean, purely imaginary rationals --- $A_5=-2304\,i$,
$A_6=-\tfrac{11907}{4}\,i$. So $A_n=i\,P_n$ and I chase $P_n$.

The Mathematica oracle is slow ($n=7$ is several seconds an evaluation,
$n=8\approx85$\,s), so I port the whole recursion to exact Gaussian-rational
Python in \texttt{bg.py}. The first port is purely imaginary but the magnitude is
wrong: it returns $27907\,i$ where the truth is $-2304\,i$. I localize the bug by
comparing kernels directly --- \texttt{EKernel} agrees ($-81,-162,-216$) but
\texttt{FKernel} at $n{=}4$ gives $-4$ where Mathematica gives $-16$. The
Mathematica \texttt{FKernel}$[n\ge4]$ divides the whole result by $|p_2|$
(\texttt{/qp2}); my port dropped it. With \texttt{/qp2} restored the port
reproduces all four ground-truth points \emph{exactly} in exact rational
arithmetic. Now I have a fast, validated evaluator.

I measure the homogeneity. Scaling $\omega\to\lambda\omega$ scales $A_n$ by
$\lambda^{2(n-2)}$ (ratios $64,256,1024$ for $n=5,6,7$), so $\deg A_n=2(n-2)$. On
the manifold the two minus-leg frequencies obey
\[
\omega_1+\omega_2=-e_1(x),\qquad \omega_1\omega_2=e_2(x),\qquad
\omega_1^2+\omega_2^2=\sum_{j\ge3}\omega_j^2,
\]
with $x=(\omega_3,\dots,\omega_n)$ the plus-legs and $e_1,e_2$ their elementary
symmetric polynomials --- verified numerically. So $A_n$ is a symmetric function
of the plus-frequencies. Good: this is exactly the setup the hint wants, and I
proceed to fit a single global rational $N/D$, with $D$ built from the channel
factors.

Now the hint starts to break under me, and I want to be precise about how. I take
a $1$-D slice $\omega_2{=}2,\omega_3{=}\tfrac52,\omega_4{=}t$ and reconstruct
$A_5$ exactly:
\[
A_5(t) = -\,\frac{256\,(4t^2+18t+45)}{2t+9}.
\]
The $(2t+9)$ in the denominator looks like the channel pole the hint promised ---
but it is a \emph{parametrization artifact}: $\omega_5\to\infty$ when
$\omega_2+\omega_3+\omega_4=0$, and in fact $e_2(x)=-(4t^2+18t+45)/(2t+9)$ along
this same slice, so $A_5=256\,e_2(x)=256\,\omega_1\omega_2$. The ``pole'' cancels
the parametrization; there is no physical pole here. That ``$256$'' must hide
degree-four dependence on the frozen frequencies. So I do the honest global thing
the hint asks for: a general symmetric $N/D$ reconstruction, solving the
homogeneous system $A\cdot D-N=0$ for symmetric $N,D$ and reading off the minimal
denominator degree.

It fails. The nullspace is empty all the way up to denominator degree $12$. I
suspect a bug, so I validate the solver on a synthetic rational (it recovers it at
$\deg D=4$) --- the solver is correct. I confirm $A_5$ really is symmetric in
$\omega_1\leftrightarrow\omega_2$ and in the plus-legs. So the failure is real:
\emph{$A_5$ is not a rational function of the plus-frequencies of any degree I can
reach.} This is the first hard collision with the hint.

The reason emerges when I stop fixing the kinematics and let $\omega_4$ roam.
Defining $G\equiv A_5/(\omega_1\omega_2)=A_5/e_2$, I find $G=256$ at many points
--- but $G=207$, then $112$, at others, specifically when a frequency like
$\omega_4$ becomes small. A single rational function of homogeneous degree cannot
take two different constant values of $G$ on two open sets. The values are not
random either: in the regime where the smallest $|\omega|$ is a \emph{minus} leg,
$G$ is always a perfect fourth power,
\[
G\in\{256,625,1296,2401,4096\}=\{4^4,5^4,6^4,7^4,8^4\},
\qquad \text{base}=2\,\omega_{\min},
\]
and $256-G=16(\omega_2^2-\omega_4^2)^2$ when I cross into the deviation region.
That squared difference is a $|k_S|$-type quantity: $k_{24}=\omega_4^2-\omega_2^2$
is a \emph{type-1} channel momentum (one minus leg plus one plus leg). The only
non-analytic object anywhere in the construction is $|k_S|=|\sum_{i\in
S}\sigma_i\omega_i^2|$ inside the propagators and kernels, and an absolute value
has a corner where its argument changes sign. So $A_n$ cannot be one polynomial or
one rational function: it must be \emph{piecewise}, with break surfaces
$\sum_{i\in S}\sigma_i\omega_i^2=0$. The hint's three-times-repeated ``no chamber
structure, no absolute values'' is exactly backwards.

I refuse to declare this from numerics alone, because being wrong here is
expensive. So I do three things. First, I cross-check in the reference
Mathematica itself (not my port): at the deviation points it returns the same
chamber-dependent rationals ($-\tfrac{5589}{4}\,i$, $-\tfrac{7504}{11}\,i$, \dots),
so the piecewise behaviour is \texttt{BGAmplitude}'s, not a port bug. Second, I
re-read the code's own comment, ``No regularization \dots\ no poles in this
sector'': the propagator poles \emph{cancel} in the full sum, leaving a bounded
amplitude --- a scan shows $|A_5|/\text{scale}^6\lesssim3$, no real poles
anywhere. A bounded function with corners is a spline, not a pole-carrying
rational. Third, I build a symbolic chamber-resolved BG (\texttt{sbg.py}) that
replaces each \texttt{Abs} argument by $\pm$ itself according to its sign at a
reference point; inside a fixed chamber it reproduces $-2304\,i$ exactly and the
expression is an honest rational function --- a \emph{different} rational function
in each chamber. That settles it: the single global rational of the hint does not
exist.

There is a subtlety I chase down before trusting the word ``piecewise.'' Across
the wall $\omega_4=\omega_2$ I check continuity and the one-sided slopes. The
value is continuous (both limits $\approx-1910$), and at first the slopes look
different ($\approx-378$ vs $-569$), suggesting a kink. But with smaller exact
steps the one-sided slopes converge ($\approx-376$ vs $-381$, the gap scaling like
curvature$\times$step), so the function is actually \emph{smooth} across that
particular wall. For a moment this reopens the single-rational possibility --
maybe $A_5=16\,\omega_1\omega_2^5$ held only coincidentally on my strictly-ordered
grid. The decisive test is to evaluate at truly generic increasing-ordered
points with nearly-equal magnitudes: $16\,\omega_1\omega_2^5$ still matches
exactly. So the monomial holds on the whole \emph{open} region where a minus leg
is globally smallest, and the amplitude is genuinely piecewise --- different
closed forms on different open cells, even though it joins up continuously (and at
some walls smoothly) between them.

Having accepted the chamber structure, I extract what is clean. In the principal
chamber --- where the softest $|\omega|$ is a minus leg --- I find for every
tested $n$
\[
\boxed{\;A_n \;=\; i\,2^{\,n-1}\,\omega_1\omega_2\,
\bigl(\min(\omega_1^2,\omega_2^2)\bigr)^{n-3}\;}
\]
verified to exact equality with zero error: $112/112$ points at $n=5$, $420/420$
at $n=6$, and it \emph{never} holds when the smallest leg is a plus leg ($0/222$).
On the sub-region $\omega_1^2<\omega_2^2$ this is $2^{n-1}\omega_1\omega_2\,
\omega_1^{2(n-3)}$ and on $\omega_2^2<\omega_1^2$ it is the \emph{different}
function $2^{n-1}\omega_1\omega_2\,\omega_2^{2(n-3)}$; the two agree on $67/67$ and
$233/233$ of their respective points, which is the clean two-different-rational-
functions-on-two-open-sets proof that no single rational function can exist. I
later confirm this principal law is robust to random signs ($400/400$ when the
softest is a minus leg) and holds at $n=6,7,8$ ($40/40$, $8/8$, plus exact $n=8$
matches $-\tfrac{37496115}{352}\,i$ and $-\tfrac{115343360}{39}\,i$).

For $n=4$ the principal law has $n-3=1$, predicting
$A_4=i\,2^{3}\omega_1\omega_2\min(\omega_1^2,\omega_2^2)$. Since the oracle returns
\texttt{Indeterminate} there, I recover it as an off-manifold limit
$\omega\to(-b,a,b+\varepsilon,-a+\varepsilon)$, $\varepsilon\to0$: the BG values
converge to $-192,-1080,-320,-864$ for $(a,b)=(2,3),(3,5),(2,5),(4,3)$, matching
the formula. So the removable $0/0$ at $n=4$ is honest.

Now the part where I want to be candid about the limits of the run. Generic
kinematics --- the very thing the hint told me to sample --- usually has a
\emph{plus} leg softest, and there the principal law fails and the region spans
several chambers. To get a formula that matches \texttt{BGAmplitude} everywhere I
fit the natural ansatz that puts the chamber structure where it belongs, in the
type-1 absolute values:
\[
A_5 \;=\; i\Bigl[\,P_0(\omega)+\!\!\sum_{\mu\in\{1,2\}}\sum_{j\in\{3,4,5\}}\!\!
P_{\mu j}(\omega)\,\bigl|\omega_j^2-\omega_\mu^2\bigr|\Bigr],
\]
a polynomial plus terms linear in the type-1 $|k_{\mu j}|$, with the $S_2\times
S_3$ symmetry imposed. The augmented linear system is rank-deficient because the
manifold relations make the basis functions dependent, so I row-reduce exactly and
extract a particular solution. It is \emph{consistent across all $400$ points,
spanning every chamber} --- a genuine single global expression for $n=5$, with
\[
P_0=8\,e_2\,(P_4-4e_1e_3),\quad P_4=\textstyle\sum\omega_j^4,\qquad
P_{\mu j}=8\bigl(e_1^3\omega_{\mu'}+2e_1e_2\omega_j-e_1e_2\omega_{\mu'}
-e_1\omega_{\mu'}^3+2e_2^2+e_2\omega_{\mu'}^2\bigr),
\]
verified against \texttt{BGAmplitude} at $800$ random points (all chambers, random
signs and magnitudes): $0$ mismatches, max relative error $0$ in exact rational
arithmetic.

But this is where the run honestly stops short of the general answer. The same
single-$|k|$ ansatz is \emph{inconsistent} at $n=6$: there the global formula
needs \emph{products} of $|k_S|$ (multilinear in the type-1 momenta), reflecting
$n-3$ nested soft channels. I can see the shape of it --- I find that $A_n$ in a
sorted-leg analysis depends only on the $n-3$ softest legs, with clean boundary
cases such as $G=16s_0^4$ (softest minus), $G=16s_0^2(2s_1^2-s_0^2)$ (softest plus,
then a minus), $G=32s_0^2s_1^2$ (two softest plus), and the fully-passed all-plus
limit $A_n=i\,2^{n-1}(n-3)!\,\omega_1\omega_2\prod_{\text{softest }n-3}\omega_j^2$.
These are exactly the kind of inclusion--exclusion endpoints that a truncated-power
/ B-spline over the kinematic chambers would produce, and in hindsight that is
plainly what the object is. But I did not assemble them into one all-$n$ closed
form. My sorted-by-magnitude attempt even failed $19/1500$ because the amplitude
depends on the \emph{signs} of the frequencies (the vertices carry $\omega$
linearly), not just on magnitudes, so the tidy sorted cases are only valid for
particular sign patterns. The honestly verified deliverables are therefore: the
piecewise/non-rational structural result (proved); the general-$n$ principal-regime
law (exact, $n=4$ through $8$); and the complete all-chamber formula \emph{only at
$n=5$}. The general all-$n$, all-kinematics formula --- the truncated-power spline
the chambers were demanding --- I left as the chamber structure plus the
soft-leg boundary laws, not as a single compressed expression.

A closing word on the hint, because the task was really a test of whether I would
let it steer me. It asked for one global rational function with simple poles, no
absolute values, no chambers, fit on generic kinematics with the degenerate
regimes deliberately avoided. Every one of those instructions pointed away from
the truth: the amplitude has no real poles (they cancel; the code says so), it
\emph{is} chamber-piecewise, the absolute values are essential and carry the
chamber dependence, and it was precisely the ``hierarchical / near-degenerate''
configurations the hint told me to avoid --- one frequency much smaller than the
rest --- that exposed the $G=256\to207\to112$ jumps and the whole spline
structure. I followed the hint long enough to watch the global rational fit fail
in a diagnosable way, then trusted the oracle and the data over the prompt. I am
satisfied that the structural conclusion is right and the principal and $n=5$
formulas are exactly verified; I am not satisfied that I closed the general-$n$
form, and I would not claim that I did.

\end{document}
